% Options for packages loaded elsewhere
\PassOptionsToPackage{unicode}{hyperref}
\PassOptionsToPackage{hyphens}{url}
\PassOptionsToPackage{dvipsnames,svgnames,x11names}{xcolor}
%
\documentclass[
  10pt,
]{article}
\usepackage{amsmath,amssymb}
\usepackage{iftex}
\ifPDFTeX
  \usepackage[T1]{fontenc}
  \usepackage[utf8]{inputenc}
  \usepackage{textcomp} % provide euro and other symbols
\else % if luatex or xetex
  \usepackage{unicode-math} % this also loads fontspec
  \defaultfontfeatures{Scale=MatchLowercase}
  \defaultfontfeatures[\rmfamily]{Ligatures=TeX,Scale=1}
\fi
\usepackage{lmodern}
\ifPDFTeX\else
  % xetex/luatex font selection
\fi
% Use upquote if available, for straight quotes in verbatim environments
\IfFileExists{upquote.sty}{\usepackage{upquote}}{}
\IfFileExists{microtype.sty}{% use microtype if available
  \usepackage[]{microtype}
  \UseMicrotypeSet[protrusion]{basicmath} % disable protrusion for tt fonts
}{}
\makeatletter
\@ifundefined{KOMAClassName}{% if non-KOMA class
  \IfFileExists{parskip.sty}{%
    \usepackage{parskip}
  }{% else
    \setlength{\parindent}{0pt}
    \setlength{\parskip}{6pt plus 2pt minus 1pt}}
}{% if KOMA class
  \KOMAoptions{parskip=half}}
\makeatother
\usepackage{xcolor}
\usepackage[margin=1in]{geometry}
\usepackage{longtable,booktabs,array}
\usepackage{calc} % for calculating minipage widths
% Correct order of tables after \paragraph or \subparagraph
\usepackage{etoolbox}
\makeatletter
\patchcmd\longtable{\par}{\if@noskipsec\mbox{}\fi\par}{}{}
\makeatother
% Allow footnotes in longtable head/foot
\IfFileExists{footnotehyper.sty}{\usepackage{footnotehyper}}{\usepackage{footnote}}
\makesavenoteenv{longtable}
\setlength{\emergencystretch}{3em} % prevent overfull lines
\providecommand{\tightlist}{%
  \setlength{\itemsep}{0pt}\setlength{\parskip}{0pt}}
\setcounter{secnumdepth}{-\maxdimen} % remove section numbering
\ifLuaTeX
  \usepackage{selnolig}  % disable illegal ligatures
\fi
\usepackage{bookmark}
\IfFileExists{xurl.sty}{\usepackage{xurl}}{} % add URL line breaks if available
\urlstyle{same}
\hypersetup{
  colorlinks=true,
  linkcolor={black},
  filecolor={Maroon},
  citecolor={Blue},
  urlcolor={black},
  pdfcreator={LaTeX via pandoc}}

\author{}
\date{}

\begin{document}

\section{BA0-T3 - Base Analysis responsibility and non-goals closure
review}\label{ba0-t3---base-analysis-responsibility-and-non-goals-closure-review}

\textbf{Revision:} R1

\textbf{Status:} CLOSED

\textbf{Repository baseline reviewed:}
\texttt{ced13f9329c89c759b3608d66c3f81ac79b9469b}

\textbf{Phase:}
\texttt{BA0\ -\ Base\ Analysis\ responsibility\ and\ boundary}

\textbf{BA1 status at completion:} \texttt{NOT\ STARTED}

\subsection{1. Purpose and closure
question}\label{purpose-and-closure-question}

BA0-T3 reviews BA0-R, BA0-T1 and BA0-T2 together and asks only whether
the Base Analysis responsibility boundary is now precise enough to begin
ontology work without importing a pre-approved type list, a
threat-method model, a tool-native representation, or a second project
source of truth.

The closure question is:

\begin{quote}
What is the minimum methodology-neutral responsibility of Base Analysis
between governed documentation, progressive human understanding,
method-specific analysis, corrective feedback and change-aware
re-analysis?
\end{quote}

This review does not design BAE types, relations, view schemas,
provenance serialization, Common Finding, STRIDE overlays or
implementation structures.

\subsection{2. Evidence basis}\label{evidence-basis}

The review uses:

\begin{itemize}
\tightlist
\item
  \href{../literature/syntheses/BA0_R_CLOSURE_R1.md}{BA0-R closure};
\item
  \href{BA0_T1_FR3_4_APPLICATION_TRIAL_R1.md}{BA0-T1 trial} and its
  provisional ledger;
\item
  \href{BA0_T2_ANALYSIS_CONSUMPTION_TEACHABILITY_GOVERNED_FEEDBACK_TRIAL_R1.md}{BA0-T2
  trial} and its provisional ledger;
\item
  \href{BA0_BASE_ANALYSIS_WORKING_HYPOTHESIS_R2.md}{BA0 Working
  Hypothesis R2};
\item
  \href{../research/research-questions.md}{research questions};
\item
  \href{../research/contributions.md}{candidate contributions};
\item
  \href{../research/hypotheses.md}{candidate hypotheses};
\item
  \href{DDTA_RESEARCH_WORK_PLAN_AFTER_DOCUMENTATION_CLOSURE_R4.md}{Work
  Plan R4}.
\end{itemize}

No Chapter 2-4 reopen criterion is triggered by this review. No new
external literature claim is introduced.

\subsection{3. Exit criteria inherited from BA0-R and
R4}\label{exit-criteria-inherited-from-ba0-r-and-r4}

BA0 can close only if the review can decide, without ontology design:

\begin{enumerate}
\def\labelenumi{\arabic{enumi}.}
\tightlist
\item
  the responsibility of Base Analysis;
\item
  the authority boundary between governed documentation and analysis;
\item
  what semantic canonicality means;
\item
  why a shared analytical representation is needed rather than
  independent re-interpretation by every consumer;
\item
  how origin/provenance and unresolved uncertainty are handled at
  responsibility level;
\item
  whether human progressive views, method projections, diagnostics,
  corrective feedback and change impact belong to the responsibility
  boundary;
\item
  explicit non-goals;
\item
  whether remaining uncertainty is properly deferred to BA1-BA6.
\end{enumerate}

\subsection{4. Falsification review}\label{falsification-review}

\subsubsection{4.1 Could Base Analysis become a second project source of
truth?}\label{could-base-analysis-become-a-second-project-source-of-truth}

Yes, if analytical normalization or tool inference were allowed to
create project commitments silently.

BA0-R and T1 already reject that move. T2 reinforces the boundary by
returning the recognition-unavailability question as a corrective
documentation candidate instead of inventing an access policy.

\textbf{Disposition:} the risk is controlled only if governed
documentation remains project authority and every Base Analysis
proposition that claims project meaning retains source/origin
provenance. This responsibility is retained.

\subsubsection{4.2 Does BA0 need one physically persistent canonical
graph?}\label{does-ba0-need-one-physically-persistent-canonical-graph}

No.~T1 shows that stable shared meaning is useful while persistence
remains unproven. T2 demonstrates that human and method-specific
projections can consume the same accepted semantics without requiring
one storage technology.

\textbf{Disposition:} semantic canonicality is retained; physical
canonical persistence is rejected as a BA0 requirement.

\subsubsection{4.3 Does every useful semantic responsibility need a
first-class BAE
type?}\label{does-every-useful-semantic-responsibility-need-a-first-class-bae-type}

No.~BA0-R and T1 directly falsify this. Transport, realization,
correlation, failure semantics, security applicability and
responsibility separation may be analytically relevant without implying
autonomous metaclasses.

\textbf{Disposition:} no BAE type is accepted by BA0-T3. Representation
form remains BA1-BA4 work.

\subsubsection{4.4 Must method-specific semantics enter the shared core
to make analysis
possible?}\label{must-method-specific-semantics-enter-the-shared-core-to-make-analysis-possible}

No.~T2's bounded STRIDE probe consumes shared
source/destination/information/correlation/responsibility semantics and
keeps STRIDE dispositions outside the common core.

\textbf{Disposition:} method neutrality and consumer isolation are
retained as Base Analysis responsibilities.

\subsubsection{4.5 Is progressive teachability itself a Base Analysis
workflow?}\label{is-progressive-teachability-itself-a-base-analysis-workflow}

Not as a visualization or usability workflow. T2 only demonstrates that
the same semantics can support an overview-to-source ladder.

However, if Base Analysis could not support bounded non-authoritative
projections and source drill-down, the thesis would require a second
independently authored model for human comprehension, undermining shared
meaning and provenance.

\textbf{Disposition:} retain only \textbf{projection readiness and
drill-down support} as a Base Analysis responsibility. Exact view types,
rendering and empirical usability remain BA4/evaluation concerns.

\subsubsection{4.6 Does Base Analysis own governed corrective
feedback?}\label{does-base-analysis-own-governed-corrective-feedback}

No.~R2 is too broad if read as assigning the correction workflow to Base
Analysis. Governed documentation already owns project decisions and
acceptance.

T2 supports a narrower responsibility: an analysis result, insufficiency
or diagnostic must be able to point back through shared semantics to the
governed source area that requires review.

\textbf{Disposition:} replace \texttt{governed\ corrective\ feedback} as
an operational BA responsibility with \textbf{source-localized feedback
handoff}. Document ownership, decision-making, acceptance and lifecycle
remain in documentation governance and later analysis workflow.

\subsubsection{4.7 Is a reviewed analytical addition a required Base
Analysis origin
class?}\label{is-a-reviewed-analytical-addition-a-required-base-analysis-origin-class}

Not proven. T1 rejected an unnecessary invented endpoint. T2 encountered
missing project semantics and correctly returned a documentation gap
rather than adding an analysis-layer project fact.

RQ1 requires explicit review of derived elements, but the current
evidence does not require a standing class of shared-core analytical
additions that introduce new project meaning.

\textbf{Disposition:} \texttt{grounded}, \texttt{derived} and
\texttt{diagnostic/unresolved} responsibilities are retained.
\texttt{reviewed\ analytical\ addition} is \textbf{DEFERRED / NOT
REQUIRED BY BA0}. If BA1-BA6 later establishes a necessary
methodology-neutral analytical construct not directly grounded in
project commitments, it must remain explicitly
non-project-authoritative, provenance-bearing and reviewed; that future
case may trigger a BA0 reopen if it changes the responsibility boundary.

\subsubsection{4.8 Must BA0 automatically detect every contradiction or
ambiguity in governed
documentation?}\label{must-ba0-automatically-detect-every-contradiction-or-ambiguity-in-governed-documentation}

No.~Reliable semantic contradiction detection over arbitrary narrative
would exceed the current thesis scope and would collapse BA into an
extraction/validation engine.

BA0 does need a rule for what happens when conflicting, ambiguous or
insufficient evidence is encountered during derivation or consumption:
it must not be silently resolved into accepted project meaning.

\textbf{Disposition:} retain \textbf{explicit unresolved/diagnostic
state and localization}. Automatic discovery of all inconsistencies is
not a BA0 requirement.

\subsubsection{4.9 Does BA0 own re-analysis orchestration after
changes?}\label{does-ba0-own-re-analysis-orchestration-after-changes}

No.~RQ4 requires that provenance and baselines identify affected
analytical artifacts; it does not require Base Analysis itself to
schedule or execute analyses.

\textbf{Disposition:} retain \textbf{change-impact traceability} as a
Base Analysis responsibility; exact identity lifecycle, stale-state
rules and re-execution orchestration remain BA3, BA6 and the later
analysis envelope.

\subsection{5. Disposition of the R2 working
invariants}\label{disposition-of-the-r2-working-invariants}

\begin{longtable}[]{@{}
  >{\raggedright\arraybackslash}p{(\columnwidth - 4\tabcolsep) * \real{0.3333}}
  >{\raggedright\arraybackslash}p{(\columnwidth - 4\tabcolsep) * \real{0.3333}}
  >{\raggedright\arraybackslash}p{(\columnwidth - 4\tabcolsep) * \real{0.3333}}@{}}
\toprule\noalign{}
\begin{minipage}[b]{\linewidth}\raggedright
R2 item
\end{minipage} & \begin{minipage}[b]{\linewidth}\raggedright
T3 disposition
\end{minipage} & \begin{minipage}[b]{\linewidth}\raggedright
Closure interpretation
\end{minipage} \\
\midrule\noalign{}
\endhead
\bottomrule\noalign{}
\endlastfoot
BA0-I1 Source authority & RETAIN & Governed documentation remains
project authority. \\
BA0-I2 Analysis-layer semantic canonicality & RETAIN / REFINE & Stable
accepted baseline-scoped identity and meaning; persistence optional. \\
BA0-I3 Explicit origin and provenance & RETAIN & Required for
drill-down, review, feedback and change impact. \\
BA0-I4 No silent inference authority & MERGE & Consequence of source
authority, provenance and explicit unresolved state. \\
BA0-I5 Method neutrality & RETAIN & Method-owned semantics stay outside
the shared core. \\
BA0-I6 Representation independence & MOVE TO CLOSED NON-GOAL & No graph
DB, notation, renderer or serialization is mandated by BA0. \\
BA0-I7 Bounded modeling scope & RETAIN AS MINIMALITY CONSTRAINT &
Include only semantics justified by shared
understanding/analysis/provenance/change needs. \\
BA0-I8 Change alignment and scoped revalidation & RETAIN / REFINE & BA
must support impact identification; it does not own re-execution
orchestration. \\
BA0-I9 Diagnostic visibility and localization & RETAIN / REFINE &
Encountered ambiguity/conflict/insufficiency must remain explicit;
universal detection is not required. \\
BA0-I10 No method/tool semantic ownership & MERGE & Consequence of
source authority, method neutrality and representation independence. \\
BA0-I11 Progressive teachability & RETAIN / REFINE & Projection
readiness and source drill-down, not a visualization/usability claim. \\
BA0-I12 Governed corrective feedback & REDUCE & BA supports
source-localized feedback handoff; governance remains outside BA. \\
BA0-I13 Consumer isolation with shared meaning & MERGE & Covered by
shared semantic identity plus method neutrality. \\
\end{longtable}

The reduction from thirteen working invariants to a smaller closed
responsibility set is intentional. BA0 closure should define the
boundary, not pre-design the metamodel.

\subsection{6. Closed Base Analysis responsibility
statement}\label{closed-base-analysis-responsibility-statement}

For a governed documentation baseline, \textbf{Base Analysis is the
accepted methodology-neutral analytical representation of the shared
project meaning required by DDTA}. It preserves baseline-scoped semantic
identity, source/origin provenance and explicit unresolved state
sufficiently to support reproducible reuse, progressive human and
method-specific projections, source drill-down, change-impact reasoning
and source-localized feedback handoff.

Base Analysis is \textbf{not project authority}. Governed documentation
is. Base Analysis may normalize and derive analytical structure from
governed meaning, but it must not silently create project commitments.
Method-specific interpretations and analysis outputs remain outside the
shared core.

The representation may be reused or reproducibly rebuilt. Semantic
canonicality does not require one permanently materialized graph, one
storage technology, one diagram or one modeling notation.

\subsection{7. Closed responsibility
set}\label{closed-responsibility-set}

\subsubsection{BA0-C1 - Authority and provenance
boundary}\label{ba0-c1---authority-and-provenance-boundary}

Project facts represented in Base Analysis must be grounded in governed
documentation, with origin and source traceability. Derived analytical
normalization must remain distinguishable from source-stated meaning. No
tool, method or analyst inference may silently become project truth.

\subsubsection{BA0-C2 - Baseline-scoped shared semantic
identity}\label{ba0-c2---baseline-scoped-shared-semantic-identity}

For a governed baseline, analytically relevant referents and
propositions must have stable enough accepted identity and meaning to be
reused across consumers, projections, comparison and revalidation. Exact
identity/equivalence mechanics remain deferred.

\subsubsection{BA0-C3 - Method-neutral shared
core}\label{ba0-c3---method-neutral-shared-core}

Base Analysis contains shared project semantics rather than STRIDE,
STRIDE-AI or other method-owned threat semantics. Consumers may project
the common meaning and attach isolated method-specific interpretation
without redefining the core.

\subsubsection{BA0-C4 - Explicit uncertainty and diagnostic
localization}\label{ba0-c4---explicit-uncertainty-and-diagnostic-localization}

When derivation or consumption encounters conflicting, ambiguous,
missing or insufficient evidence, Base Analysis must preserve that
unresolved status and enough provenance/context to localize review
rather than silently selecting a project fact. BA0 does not require
universal automatic inconsistency detection.

\subsubsection{BA0-C5 - Projection readiness and source
drill-down}\label{ba0-c5---projection-readiness-and-source-drill-down}

The same accepted semantics must support bounded, non-authoritative
projections for progressive human understanding and method consumption,
with a path back to exact governed sources. Exact view contracts and
usability claims remain deferred.

\subsubsection{BA0-C6 - Change-impact
traceability}\label{ba0-c6---change-impact-traceability}

The representation must preserve enough source, semantic and dependency
linkage to identify which analytical meaning may require review after a
relevant governed-source change and to support later targeted
re-analysis. BA0 does not define the material stale-state/re-execution
mechanism.

\subsubsection{BA0-C7 - Source-localized feedback
handoff}\label{ba0-c7---source-localized-feedback-handoff}

Analysis results, insufficiencies and diagnostics must be able to
reference the shared semantics and governed source area that motivates
review. Base Analysis does not choose the corrective policy, document
kind, acceptance decision or governance lifecycle.

\subsubsection{BA0-C8 - Minimality and representation
independence}\label{ba0-c8---minimality-and-representation-independence}

Base Analysis includes only methodology-neutral semantics justified by
shared project understanding, multi-consumer analysis, provenance,
diagnostics, projection or change/re-analysis needs. It is not required
to be a complete systems model, and its responsibility is independent of
one storage, graph, notation, renderer or implementation technology.

\subsection{8. Origin-state closure}\label{origin-state-closure}

At BA0 responsibility level:

\begin{itemize}
\tightlist
\item
  \textbf{grounded:} REQUIRED;
\item
  \textbf{derived:} REQUIRED, with explicit provenance and no new
  project commitment;
\item
  \textbf{diagnostic / unresolved:} REQUIRED as an explicit state when
  accepted meaning cannot be justified;
\item
  \textbf{reviewed analytical addition:} DEFERRED / NOT REQUIRED BY BA0.
\end{itemize}

This is not a material schema. BA3 must define the concrete
provenance/review model. BA1 must not interpret these labels as BAE
metaclasses by default.

\subsection{9. Closed non-goals}\label{closed-non-goals}

Base Analysis is not responsible for:

\begin{itemize}
\tightlist
\item
  replacing governed project documentation or its governance workflow;
\item
  reliable automatic extraction or migration from arbitrary, legacy or
  unstructured narrative;
\item
  silently accepting LLM/NLP/tool/analyst inference as project truth;
\item
  automatically discovering every contradiction or ambiguity in
  arbitrary text;
\item
  defining the final BAE type taxonomy in BA0;
\item
  requiring every useful semantic responsibility to become a first-class
  metaclass;
\item
  copying SysML, KerML, ArchiMate, AADL or another general-purpose
  modeling language;
\item
  being a STRIDE DFD, STRIDE-AI model or universal threat-analysis
  ontology;
\item
  defining AnalysisRecord, Common Finding,
  Finding-to-SecurityRequirement acceptance or formal overlay schemas;
\item
  governing corrective project decisions or choosing the target document
  type;
\item
  proving universal support for every threat-analysis method;
\item
  enforcing one persisted graph, database, serialization, diagram or
  renderer;
\item
  empirically proving that a particular visualization reduces reading
  time or improves comprehension;
\item
  executing or scheduling downstream re-analysis itself.
\end{itemize}

\subsection{10. Questions deliberately deferred after
BA0}\label{questions-deliberately-deferred-after-ba0}

The following are not BA0 closure blockers:

\begin{itemize}
\tightlist
\item
  which responsibilities become first-class BAE types, relations, roles,
  properties, constraints or projections - BA1;
\item
  minimal relation and canonical action vocabulary - BA2;
\item
  exact identity/equivalence/lifecycle rules, source locators,
  provenance materialization, review state and stale state - BA3;
\item
  exact progressive-human and method-specific view/projection contracts
  - BA4;
\item
  controlled vocabulary and optional authoring/extraction assistance -
  BA5;
\item
  holdout regression, broader method-neutrality pressure and full Base
  Analysis closure - BA6;
\item
  AnalysisRecord, Common Finding and formal overlay contracts - after
  BA6.
\end{itemize}

\subsection{11. Alignment with research questions and
hypotheses}\label{alignment-with-research-questions-and-hypotheses}

The closed boundary preserves the thesis commitments without expanding
them:

\begin{itemize}
\tightlist
\item
  RQ1/H1: governed documentation can yield a reviewed
  methodology-neutral Base Analysis with provenance and explicit review
  of derived structure;
\item
  RQ2/H2: distinct methods can consume the same semantic core while
  method-owned meaning remains isolated;
\item
  RQ3/H3: later accepted findings can retain provenance through Base
  Analysis to governed project subjects without making BA the governance
  owner;
\item
  RQ4/H4: baseline/source/dependency traceability can support
  identification of impacted analytical artifacts after change.
\end{itemize}

The closure does not claim that these research questions or hypotheses
are already empirically proven.

\subsection{12. Closure decision}\label{closure-decision}

The falsification review found two R2 overextensions - operational
ownership of governed corrective feedback and the implied necessity of a
reviewed analytical-addition class - but neither requires a new BA0
experiment once narrowed as above.

The remaining open questions concern representation choices, schemas,
projection contracts and regression evidence rather than the
responsibility boundary itself.

\textbf{Decision:
\texttt{BA0\ RESPONSIBILITY\ AND\ NON-GOALS\ =\ CLOSED}.}

\textbf{No BAE type or relation is accepted by BA0-T3.}

\textbf{BA1 remains \texttt{NOT\ STARTED} by this microstep.}

The next authorized microstep after this closure is:

\begin{quote}
\textbf{BA1-T1 - minimal BAE ontology candidate derivation from the
closed BA0 responsibilities.}
\end{quote}

BA1-T1 must begin from BA0-C1..C8, BA0-R prior-art pressure and the
governed DDTA corpora. Historical
\texttt{actor/component/asset/boundary/data\_flow} categories,
ThreatForge structures, STRIDE DFD elements and systems-modeling
metaclasses are candidates only if independently justified; none is
inherited as accepted ontology.

\subsection{13. Reopen criteria}\label{reopen-criteria}

Reopen BA0 only if a later phase produces a material
responsibility-level counterexample, for example:

\begin{enumerate}
\def\labelenumi{\arabic{enumi}.}
\tightlist
\item
  BA1-BA3 requires a methodology-neutral semantic responsibility that
  cannot fit BA0-C1..C8 without changing the boundary;
\item
  a concrete method consumer can operate only by moving method-owned
  semantics into the shared core;
\item
  provenance/identity work shows that governed-document authority cannot
  be preserved under the closed boundary;
\item
  BA4 cannot generate coherent human/method projections without
  introducing a second independently authored project model;
\item
  BA6 holdout regression reveals that the closed responsibilities are
  insufficient or mutually inconsistent.
\end{enumerate}

Ordinary choices among types, relations, properties, storage
technologies or renderings do not reopen BA0.

\end{document}
